% Explain how the raw data was loaded, cleaned and transformed.

\subsection{The Road Traffic Fine Management Dataset}

The project uses the Road Traffic Fine Management Process event log published by \textcite{deleoni2015}. The dataset records real process executions of traffic fine cases handled by an Italian local police authority. Each case corresponds to one process instance — a trace — and each trace consists of a sequence of events representing activities such as \textit{Create Fine}, \textit{Send Fine}, \textit{Insert Fine Notification}, \textit{Add penalty}, \textit{Payment}, and \textit{Send for Credit Collection}.

The event log was provided in XES format. During data preparation, the file was loaded with \texttt{pm4py} and converted into a pandas dataframe for downstream analysis. The raw event log contains:

\begin{itemize}
    \item 561,470 events across 150,370 cases
    \item 11 distinct activity types
    \item A timeframe spanning from 2000 to 2013
\end{itemize}

\subsection{Data Cleaning and Preprocessing}

The cleaning pipeline applies the following steps in order:

\textbf{Sorting.} Events were sorted by case identifier and timestamp to preserve the chronological order within each trace. This ordering is essential because the activity sequence defines the observable process behavior.

\textbf{Deduplication.} Exact duplicate events — rows sharing the same case, activity, and timestamp — were removed.

\textbf{Column removal.} Two attributes were dropped after profiling revealed them to be uninformative: \texttt{org:resource} had a fill rate below one percent with only dummy placeholder values, and \texttt{matricola} showed a near-identical pattern. Retaining these columns would introduce noise without contributing predictive signal.

\textbf{Numeric coercion.} The \texttt{amount} field was cast to a numeric type, with unparseable entries set to zero.

\textbf{Case start time.} A derived variable \texttt{case\_start\_time} records the first timestamp of each case. This variable enables a temporal train-test split in which older cases are used for training and newer cases for evaluation, preventing information from the future from leaking into the training set.

\subsection{Attribute Selection for Modeling}

Not all columns in the event log carry useful information for prediction. After profiling each attribute for fill rate, cardinality, and predictive relevance, two groups of features were retained for the data-aware model variant:

\begin{itemize}
    \item \textbf{Numeric:} \texttt{amount} and \texttt{points} — aggregated over the known prefix of each case.
    \item \textbf{Categorical:} \texttt{vehicleClass} and \texttt{article} — encoded via one-hot encoding, with rare article categories grouped into an \textit{Other} bucket to limit dimensionality.
\end{itemize}

Several attributes were explicitly excluded to prevent data leakage or because they lacked predictive value: \texttt{totalPaymentAmount} and \texttt{paymentAmount} reveal the outcome directly, \texttt{expense} accumulates over the full case lifetime, and remaining columns such as \texttt{notificationType}, \texttt{dismissal}, and \texttt{lastSent} are near-constant or predominantly missing.

\subsection{Outcome Definition and Labeling}

The primary predictive task is binary classification: predicting whether a traffic fine case will escalate to credit collection. The labeling logic applies a strict priority rule:

\begin{enumerate}
    \item If \textit{Send for Credit Collection} appears anywhere in the trace, the case is labeled as class~1 — regardless of whether a payment also occurred. The rationale is that once debt collectors are involved, the process has deviated from the desired path.
    \item Otherwise, if \textit{Payment} appears and \textit{Send for Credit Collection} does not, the case is labeled as class~0.
    \item Cases containing neither activity represent open or unresolved instances and are excluded from the supervised learning dataset.
\end{enumerate}

After applying this labeling scheme, 126,602 completed cases remain: approximately 60\% belong to class~0 and 40\% to class~1. Since the class distribution is reasonably balanced, no resampling techniques such as SMOTE are applied. The classifiers use built-in class weighting to account for the mild skew without distorting process-level patterns.

The positive class represents escalation to credit collection. From a business perspective, failing to identify these cases early is costly, making recall for class~1 a key evaluation metric alongside overall accuracy.